% Sec. III (Symbolic Model and Notation) and Sec. IV (Formal Modelling Methodology).
\section{Symbolic Model and Notation}\label{sec:symbolic}

We model the protocol using \tamarin~\cite{meier2013}. It is a symbolic tool for
security-protocol analysis, which tracks a \emph{system} represented as a multiset of facts,
referred to as the state~$S$. The state evolves following user-defined transition \emph{rules}
of the following shape:
\begin{equation}\label{eq:rule}
  [\,L\,]\;\rule[0.55ex]{1.6em}{0.4pt}\,[\,A\,]\,\longrightarrow\;[\,R\,]
\end{equation}
with the premises~$L$, the action facts~$A$ and the conclusions~$R$.

A rule of shape~\eqref{eq:rule} can be applied when $L$ unifies with facts in~$S$. It captures
that $L$ is rewritten as $R$ in~$S$, and the action facts~$A$ are appended to the trace.
\tamarin{} features linear and persistent facts; the former are consumed upon rule application,
while the latter (prefixed with \texttt{!}) can be reused arbitrarily many times. The tool also
provides the built-in facts \code{In}, \code{Out} and \code{Fr}, capturing respectively
receiving from and sending to a public (i.e., attacker-controlled) channel, and generating fresh
and unique values (e.g., for nonces).

We further use the hashing, signing and natural-numbers built-ins. Hashing provides~$h$, used
for payment hashes $y = h(x)$ and for the commitment $h(s)$ to a revocation secret~$s$. Signing
provides \code{sign}, \code{verify} and \code{pk}, used for the channel handshake, the
commitments and the invoices. Natural numbers provide the sort of amounts (variables prefixed
with \texttt{\%}) with addition, used for the per-hop constraint between incoming amount,
outgoing amount and fee (Sec.~\ref{sec:core-amounts}).

An execution produces a trace of action facts. Security goals are then first-order formulas
over these traces (with time points \texttt{\#i}), enabling security queries over all traces
(\code{all-traces} lemmas) and reachability queries, ensuring that a certain trace exists
(\code{exists-trace} lemmas). A \emph{restriction} is a formula of the same kind that
\tamarin{} imposes on the model: only traces that satisfy every restriction are considered.

\section{Formal Modelling Methodology}\label{sec:method}

We model the protocol as a single \tamarin{} theory, \code{lightning.spthy} (37 rules), with a
single adversary. It contains both the payment-channel protocol (Sec.~\ref{sec:channel}) and
multi-hop routing (Sec.~\ref{sec:routing}). One generic rule, \code{Forward_HTLC}, handles every
forwarding hop, so routes of any length are covered. The routing layer does not abstract the
channel layer away, but the coupling is deliberately thin: commitment transactions carry no HTLC
outputs, and, besides the long-term keys, the layers interact only through three explicit
interfaces.

\paragraph{Topology and Slot} \code{Lock_Funds_And_Open} mints the persistent facts
\code{!ChannelConnect}, one per direction, which \code{Sender_Offer_HTLC}, \code{Forward_HTLC}
and \code{Open_Timeout} require; HTLCs are therefore added only on opened channels
(\code{HTLC_On_Opened_Channel}). Opening also mints one linear fact \code{Mutex(ptr)}: adding an
HTLC consumes it, and redeeming or refunding that HTLC returns it. A channel thus holds at most
one in-flight HTLC at a time and can be reused afterwards.

\paragraph{Resolution} Every rule that redeems or refunds an HTLC records the action fact
\code{Old}, and each HTLC is resolved at most once:
\begin{spthy}
restriction OneResolutionPerHTLC:
    "All ptr y #i #j. Old(ptr, y) @ #i & Old(ptr, y) @ #j
     ==> #i = #j"
\end{spthy}
This is the interface assumption between the layers. In the protocol, it holds because an HTLC
output of the current commitment can be spent only once on chain; spending it from an older
state requires publishing a revoked commitment, which the counterparty can punish
(\code{Publish_Revoked_A/B}, \code{Punish_A/B}). This argument is not mechanised, and we claim
no composition theorem. All other abstractions are also explicit, named restrictions
(Sec.~\ref{sec:restrictions}).

\subsection{Threat Model}\label{sec:threat}

\tamarin{} assumes the standard Dolev-Yao adversary model~\cite{dolevyao1983}: it controls the
whole network; it can read, modify and inject messages of arbitrary complexity; it is, however,
constrained to the term algebra under consideration. Our adversary keeps this full control.
Public keys, channel proposals and acknowledgements, and invoices are sent over the network. The
adversary also chooses the values that honest rules read from it: the balances of a channel
update (\code{Update_State}) and, at each hop, the incoming amount, the fee and the outgoing
amount (\code{Forward_HTLC}), subject only to the incoming amount being the outgoing amount plus
the fee. HTLC offers and fulfilments between adjacent hops are linear facts (\code{Route},
\code{Fulfilled}), that is, authenticated, non-replayable messages between channel partners.

\paragraph{Key Compromise} Rule \code{Compromise_Ltk} reveals the long-term key of any party at
any time (action fact \code{LtkCompromised}); the adversary can then sign as that party, e.g.,
forge its invoices: a payment for which the receiver created no invoice can then settle
(\code{Forged_Invoice_Requires_Key_Compromise}). The preimage~$x$ is never
sent on the network, and the adversary never learns it (\code{Preimage_Never_Known}). What
compromised parties can do with it is given explicitly by rule \code{AttackerSkip}: two
compromised parties that offered or forwarded HTLCs with the same payment hash may move the
fulfilment of one party's outgoing HTLC to the other party's outgoing HTLC, relaying~$x$ around
the parties between them: the wormhole attack~\cite{malavolta2019} (Sec.~\ref{sec:wormhole}).
Apart from their keys, \code{AttackerSkip} and the exemptions of Sec.~\ref{sec:restrictions},
compromised parties follow the protocol rules.

\paragraph{Time} Time is abstract: the timeout of channel \code{ptr} is the event
\code{TimedOut(ptr)} of rule \code{Open_Timeout}, which can happen at any point after the
channel is opened, and at most once per channel (\code{OneTimeoutPerPtr}). The adversary
therefore chooses when timeouts happen relative to honest actions.

\paragraph{Liveness} As mentioned before, the security of the Lightning Network relies on the
liveness of the participants, whether to deter cheating behaviour (Sec.~\ref{par:punishment}) or
to claim HTLCs (Remark~\ref{rmk:claiming}). However, a standard Dolev-Yao adversary has full
control over the network. It can notably keep a party from seeing the network (e.g., preventing
it from noticing a fraud) or from sending a message (e.g., the defending transaction), leading to
trivial attacks on the Lightning Network~\cite{mccorry2019}. Our model therefore uses \tamarin{}
restrictions to impose liveness on honest forwarders, as explained next.

\subsection{Liveness Restrictions}\label{sec:restrictions}

To make the model relevant, we restrict the adversary's control over the order of events to
allow for liveness. We also rule out pathological honest behaviour. Concretely, we impose the
following three restrictions, which constrain the routing layer (shown without the flags
\code{NO_RBT}, \code{NO_OTF}, \code{NO_IMC} that remove them for the necessity runs).
\begin{spthy}[linerange={1-3,9-14,20-26}]
restriction RedeemBeforeTimeout:
    "All ptr y #i #j.
        Redeemed(ptr, y) @ #i & TimedOut(ptr) @ #j ==> #i < #j"
#endif

#ifdef NO_OTF
// restriction removed for the necessity check
#else
restriction OutgoingTimesOutFirst:
    "All P Prev Next ptrIn ptrOut y #f #k.
        HTLCForwarded(P, Prev, Next, ptrIn, ptrOut, y) @ #f
      & TimedOut(ptrIn) @ #k
      & not (Ex #c. LtkCompromised(P) @ #c & #c < #k)
     ==> Ex #m. TimedOut(ptrOut) @ #m & #m < #k"
#endif

#ifdef NO_IMC
// restriction removed for the necessity check
#else
restriction IntermediaryMustClaim:
    "All P Prev Next ptrIn ptrOut y #f #i #k.
        HTLCForwarded(P, Prev, Next, ptrIn, ptrOut, y) @ #f
      & Redeemed(ptrOut, y) @ #i
      & TimedOut(ptrIn)  @ #k
      & not (Ex #c. LtkCompromised(P) @ #c & #c < #k)
     ==> Ex #l. Redeemed(ptrIn, y) @ #l"
\end{spthy}

\noindent\textbf{Restriction T1} (\code{RedeemBeforeTimeout}). \emph{If an HTLC is redeemed and
its channel times out, the redemption comes first.} This is not a liveness assumption but a
well-formedness condition on timeouts, and it applies to every HTLC.

\noindent\textbf{Restriction T2} (\code{OutgoingTimesOutFirst}). \emph{If the incoming HTLC of a
forwarding hop times out and the forwarder was not compromised before, then its outgoing HTLC
timed out earlier.} This is the decreasing-deadline (CLTV-delta) invariant of
Remark~\ref{rmk:deadlines}.

\noindent\textbf{Restriction T3} (\code{IntermediaryMustClaim}). \emph{If a forwarder's outgoing
HTLC is redeemed and its incoming HTLC times out, and the forwarder was not compromised before
that timeout, then its incoming HTLC is redeemed too.} This is the claiming liveness of
Remarks~\ref{rmk:claiming} and~\ref{rmk:deadlines}, on which intermediary safety rests. It only
constrains traces in which the incoming HTLC times out, the only ones in which the upstream
party can be refunded (\code{Refund_Requires_Timeout}); by T1, the claim then precedes the
timeout.

T2 and T3 do not constrain forwarders compromised before their incoming timeout. The order of a
hop's two timeouts is not assumed separately: \tamarin{} proves it from T2 and
\code{OneTimeoutPerPtr}.
\begin{spthy}
lemma ForwardTimeGap [reuse]:
    "All P Prev Next ptrIn ptrOut y #f #j #k.
        HTLCForwarded(P, Prev, Next, ptrIn, ptrOut, y) @ #f
      & TimedOut(ptrOut) @ #j
      & TimedOut(ptrIn)  @ #k
      & not (Ex #c. LtkCompromised(P) @ #c & #c < #k)
     ==> #j < #k"
\end{spthy}
As restrictions, T1 to T3 remove every trace that violates them, including traces in which an
uncompromised forwarder could not have acted in time; our results say nothing about such traces
(Sec.~\ref{par:liveness-scope}). So, in our model, the liveness of punishment
(Sec.~\ref{par:punishment}) is not encoded: no restriction mentions cheating or punishment.

The six remaining restrictions are structural. \code{Equality} and \code{Inequality} turn the
action facts \code{Eq} and \code{Neq} into constraints, which encodes checks such as signature
verification; \code{OneChannelPerPtr} and \code{OneTimeoutPerPtr} let a channel open and time
out at most once; \code{OneHTLCPerHashPerChannel} lets a payment hash be added at most once per
channel; and \code{OneResolutionPerHTLC} is the interface assumption above.

Sec.~\ref{sec:timing} removes one liveness restriction at a time. T3 is necessary for
intermediary safety: without it, \tamarin{} finds a trace with no key compromise in which an
intermediary's outgoing HTLC is redeemed while its incoming one is refunded. T1 is not needed by
any of the five liveness lemmas checked there, and T2 is not needed for intermediary safety. The
bypass theorem for routes of any length and a single offer per payment hash
(Sec.~\ref{sec:bypass-any}) is proved without T1, T2 and T3; all other results of the main theory are
proved under all restrictions.

\subsection{Structural Limitations}\label{sec:structure}

The Lightning Network is a complex protocol. Verifying it as a single monolithic block quickly
hits scalability problems with verification tools~\cite{grundmann2026}, and \tamarin{} is no
exception. This is not an assumption but an observation. When different built-ins are combined
in \tamarin{}, the search space explodes: with signatures, an equational theory,
natural-number induction and an unbounded clock, \tamarin{} exhausts memory; without the
unbounded clock, the remaining ones are tractable together. We therefore keep the block clock
out of the main theory, which represents timeouts as untimed events ordered by the restrictions above. The block clock and the CLTV arithmetic are analysed
in small auxiliary models that use hashing and natural numbers only.

Our artifact thus consists of one main theory and five small auxiliary models, all in
\code{lightning_aux.spthy}, each selected by a preprocessor flag: the PTLC countermeasure
(\code{-D=PTLC}) and its HTLC baseline (\code{-D=HTLC}), a block-clock model of CLTV deadlines
(\code{-D=CLOCK}), CLTV-delta arithmetic (\code{-D=CLTV}) and the early-timeout race
(\code{-D=TIMEOUT}). Each is self-contained and is not composed with the main theory
(Table~\ref{tab:results}). Inside the main theory, HTLC state is deliberately not folded into
the commitment transactions: \code{StateA} and \code{StateB} carry balances and revocation data
but no HTLC outputs. This keeps the channel state independent of the payments routed over it.
The price is \code{OneResolutionPerHTLC}, imposed rather than derived from the channel
mechanics; a mechanised discharge, by composition or by refinement as in~\cite{grundmann2026},
is the natural next step.

Five aspects are deliberate non-goals, each of which buys generality or tractability.
\emph{Onion privacy:} each forwarder may pick any channel partner as its next hop and nothing is
encrypted in layers, so the results cover every route but claim no privacy. \emph{On-chain
fees:} publishing a transaction is free, which separates the protocol logic from fee economics.
\emph{Fee policy:} the adversary chooses every routing fee, so the results hold for any policy.
\emph{Wallet capacity:} HTLC amounts are not checked against the balances, so the results hold
for all amounts. \emph{Per-HTLC CSV stage:} an HTLC is redeemed or refunded in one step, without
a second-stage transaction.
